%%=============================================================================
%% Inleiding
%%=============================================================================

\chapter{\IfLanguageName{dutch}{Inleiding}{Introduction}}%
\label{ch:inleiding}

Artistieke gymnastiek is een technisch veeleisende sport waarbij coördinatie, kracht, mobiliteit en timing goed op elkaar afgestemd moeten worden. Zelfs eenvoudige bewegingen vereisen een correcte uitvoering om efficiënt en veilig uitgevoerd te worden. Binnen recreatieve contexten worden gymnasten regelmatig geconfronteerd met complexe elementen waarvoor zij nog onvoldoende lichaamscontrole of kracht hebben ontwikkeld. Hierdoor ontstaan vaak technische uitvoeringsfouten.

De kip aan barre is een belangrijk overgangselement binnen de damesgymnastiek. Het element vereist een combinatie van zwaaimomentum, heupsluiting, armstrekking en actieve schouderwerking om vanuit hangpositie tot steun op de barre te komen. Vooral bij recreatieve turnsters blijkt deze beweging technisch uitdagend. Dit leidt regelmatig tot fouten zoals gebogen armen, onvoldoende heupsluiting of onvoldoende schouderdruk.

Een correcte technische uitvoering is niet alleen belangrijk voor prestatieverbetering, maar ook voor blessurepreventie. Herhaaldelijk uitvoeren van foutieve bewegingspatronen kan bijdragen aan overbelastingsletsels, vooral ter hoogte van schouder en pols bij bewegingen aan de barre. Het vroegtijdig detecteren en corrigeren van technische fouten is daarom belangrijk binnen een veilige trainingsomgeving.

Binnen recreatieve gymnastiekclubs worden technische elementen meestal in groepsverband aangeleerd. Trainers begeleiden meerdere sporters tegelijk, waardoor er beperkte tijd is voor een grondige individuele analyse van elke uitvoering. Feedback gebeurt voornamelijk op basis van visuele observatie en ervaring van de coach. Geavanceerde 3D bewegingsanalysesystemen bestaan, maar vereisen gespecialiseerde apparatuur en zijn meestal financieel niet haalbaar voor recreatieve clubs. Er is dus een duidelijke kloof tussen laboratoriumgebaseerde bewegingsanalyse en de praktische realiteit in de trainingszaal.

Recente ontwikkelingen in markerloze pose-estimatie maken het mogelijk om gewrichtsposities te bepalen op basis van standaard videobeelden, zonder gebruik van sensoren of markers. Daarnaast kunnen Large Language Models (LLM’s) informatie omzetten in begrijpelijke tekstuele feedback. De combinatie van deze technologieën biedt mogelijkheden voor het ontwikkelen van een laagdrempelig AI-gebaseerd ondersteuningssysteem voor technische analyse in recreatieve gymnastiek.

Deze bachelorproef onderzoekt in welke mate een dergelijk systeem haalbaar en bruikbaar is voor het genereren van technische feedback bij de kip aan barre.

\section{\IfLanguageName{dutch}{Probleemstelling}{Problem Statement}}%
\label{sec:probleemstelling}
Binnen recreatieve en laagcompetitieve gymnastiekclubs worden technische elementen zoals de kip aan barre aangeleerd aan turnsters van ongeveer 10 jaar en ouder. In deze context beschikken clubs doorgaans niet over gespecialiseerde bewegingsanalyseapparatuur en zijn trainers aangewezen op visuele observatie.

De doelgroep van dit onderzoek bestaat uit:

\begin{itemize}
    \item Recreatieve en laagcompetitieve turnsters
    \item Trainers binnen recreatieve gymnastiekclubs
    \item Gymnastiekclubs zonder toegang tot gespecialiseerde bewegingsanalyse
\end{itemize}

Voor deze doelgroep kan een laagdrempelig video-gebaseerd analysesysteem een meerwaarde bieden door:

\begin{itemize}
    \item Consistente detectie van vooraf gedefinieerde technische fouten
    \item Ondersteuning van trainers bij het formuleren van gerichte feedback
    \item Verhoogd bewustzijn van uitvoeringskwaliteit bij sporters
\end{itemize}

De centrale uitdaging bestaat erin te onderzoeken of een AI-gebaseerde oplossing op basis van standaard videobeelden een haalbare en bruikbare ondersteuning kan vormen binnen deze recreatieve trainingscontext.

\section{\IfLanguageName{dutch}{Onderzoeksvraag}{Research question}}%
\label{sec:onderzoeksvraag}
De centrale onderzoeksvraag luidt als volgt:

\textit{In welke mate kan een AI-gebaseerd systeem op basis van markerloze video-analyse technisch bruikbare feedback genereren voor recreatieve turnsters bij de kip aan barre?}

 Om deze onderzoeksvraag te beantwoorden worden volgende deelvragen onderzocht. Deze hebben enerzijds betrekking tot het probleemdomein en anderzijds tot het oplossingsdomein.
\newline
\textbf{Probleemdomein}:

\begin{enumerate}
    \item Welke technische fouten komen het vaakst voor bij de kip?
    \item Waarom worden terugkerende technische fouten door coaches soms niet opgemerkt, en hoe kan dit systematisch worden vastgelegd?
    \item Welke kenmerken maken technische feedback effectief en bruikbaar voor recreatieve turnsters?
\end{enumerate}

\textbf{Oplossingsdomein}:

\begin{enumerate}
    \item Welke markerloze pose-estimatie methode is het meest geschikt voor het analyseren van de kip aan barre binnen een recreatieve trainingscontext?
    \item In welke mate kunnen technische fouten automatisch worden gedetecteerd op basis van pose-data?
    \item In hoeverre kan een Large Language Model technisch correcte en bruikbare feedback geven op basis van gedetecteerde fouten?
    \item Wat zijn de beperkingen van een AI-gebaseerd feedbacksysteem binnen deze toepassing?
\end{enumerate}

\section{\IfLanguageName{dutch}{Onderzoeksdoelstelling}{Research objective}}%
\label{sec:onderzoeksdoelstelling}
Het doel van deze bachelorproef is het ontwikkelen en evalueren van een proof-of-concept van een AI-gebaseerd systeem dat:

\begin{enumerate}
    \item Gewrichtsposities extraheert uit 2D videobeelden via markerloze pose-estimatie.
    \item Vier frequent voorkomende technische fouten bij de kip aan barre rule-based detecteert:
    \begin{itemize}
        \item Gebogen armen tijdens de trekfase
        \item Onvoldoende heupsluiting
        \item Te vroege heupopening
        \item Onvoldoende schouderdruk in steunfase
    \end{itemize}
    \item Gedetecteerde fouten omzet in gestructureerde parameters.
    \item Op basis van deze parameters via een Large Language Model begrijpelijke en gerichte technische feedback genereert.
\end{enumerate}

Daarnaast wordt verkennend onderzocht in welke mate de geselecteerde fouttypes detecteerbaar zijn via 2D pose-estimatie en of de gegenereerde feedback inhoudelijk aansluit bij de gedetecteerde fouten.

Het onderzoek richt zich op de technische haalbaarheid en praktische bruikbaarheid binnen een recreatieve context, en niet op de ontwikkeling van een volledig gevalideerd diagnostisch systeem.

\section{\IfLanguageName{dutch}{Opzet van deze bachelorproef}{Structure of this bachelor thesis}}%
\label{sec:opzet-bachelorproef}

% Het is gebruikelijk aan het einde van de inleiding een overzicht te
% geven van de opbouw van de rest van de tekst. Deze sectie bevat al een aanzet
% die je kan aanvullen/aanpassen in functie van je eigen tekst.

De rest van deze bachelorproef is als volgt opgebouwd:

In Hoofdstuk~\ref{ch:stand-van-zaken} wordt een overzicht gegeven van de stand van zaken binnen het onderzoeksdomein, op basis van een literatuurstudie.

In Hoofdstuk~\ref{ch:methodologie} wordt de methodologie toegelicht en worden de gebruikte onderzoekstechnieken besproken om een antwoord te kunnen formuleren op de onderzoeksvragen.

% TODO: Vul hier aan voor je eigen hoofstukken, één of twee zinnen per hoofdstuk
In Hoofdstuk~\ref{ch:resultaten} worden de implementatie van de proof-of-concept, de regelgebaseerde foutdetectie en de feedbackgeneratie besproken.

In Hoofdstuk~\ref{ch:conclusie}, tenslotte, wordt de conclusie gegeven en een antwoord geformuleerd op de onderzoeksvragen. Daarbij wordt ook een aanzet gegeven voor toekomstig onderzoek binnen dit domein.